\documentclass[../main.tex]{subfiles}
\begin{document}
%==========================================TIÊU ĐỀ TẬP========================================
\begin{table}[h]
    \begin{adjustwidth}{-1cm}{}
        \begin{tabular}{>{\centering\arraybackslash}p{6cm}|>{\raggedright\arraybackslash}p{12cm}}
            \multirow{5}{6cm}
            % Cột bên trái: ÉPISODE và số tập
            {\\[-40pt]
                \humanist\fontsize{20pt}{20pt}\selectfont \centering
                \color{eptype}FOLGE\color{black}\\
                \houseterror\fontsize{72pt}{72pt}\selectfont
                \color{epnum}1\color{black}\\[6pt]
                \cafeteria\fontsize{20pt}{20pt}\selectfont\color{epnum}(M1-1)\color{black}
                \\[18pt]
                \humanist\fontsize{14pt}{14pt}\selectfont
                \color{epattr}WILLKOMMEN, \uline{\color{Robel} Roberu}, \uline{\color{Oga} Oga} und \uline{\color{Aruran} Aruran}!
                % \humanist\fontsize{18pt}{18pt}\selectfont
                % \color{epattr}ÉPISODE PERDU
                \color{black}
            }
             &
            % Dòng thứ nhất: tên tập tiếng Nhật
            {\fotskip\fontsize{18pt}{18pt}\selectfont この\ruby{人}{ひと}がモテる\ruby{方}{ほう}\ruby{法}{ほう}３\ruby{選}{せん}}                                                                                                                                                                         \\[6pt]
            % Dòng thứ hai: tên tập tiếng Anh
             & {\cafeteria\fontsize{18pt}{18pt}\selectfont Popularity}                                                                                                                                                                               \\[4pt]
            % Dòng thứ ba: tên tập tiếng Việt
             & {\vnmsans\fontsize{18pt}{18pt}\selectfont Nổi tiếng}                                                                                                                                                                         \\[8pt]
            % Dòng thứ tư: Nhân vật xuất hiện
             & {\caviar{\fontsize{14pt}{14pt}\selectfont Nhân vật: \emp{kaigou2} \emp{game} \emp{aruran1} \emp{roberu} \emp{oga}}} \\[8pt]
            % Dòng cuối cùng: Thời gian diễn ra của tập (thường sẽ là ngày phát hành)
             & {\continuum\fontsize{16pt}{16pt}\selectfont Ngày phát hành: 28/01/2022}                                     \\[18pt]
        \end{tabular}\\[8pt]
    \end{adjustwidth}
\end{table}
%=========================================TRUYỆN CHÍNH========================================
\color{black}

\txtbx[2]{
    {\color{vol} \uline{Tóm tắt:}} Ngày làm việc chính thức đầu tiên của Kaigou tại \hls\ biến thành buổi ``huấn luyện nổi tiếng'' cho Roberu: từ giọng nói quyến rũ, phong thái quý tộc cho đến sự tự tin. Những lời khuyên ``có một không hai'' từ Oga và Aruran chỉ khiến mọi chuyện thêm hài hước, và buổi stream thử nghiệm kết thúc bằng thảm họa. Cuối ngày, Kaigou chia sẻ về ``người đàn ông lý tưởng'' của mình, khiến ba chàng trai tò mò không dứt.
}

Chưa đầy hai ngày sau sự cố ``quả đá khổng lồ'' càn quét hành lang \hll, không khí tại căn mancave ngoại ô của \hls đã lấy lại sự yên bình hiếm hoi. Đây là ngày làm việc chính thức đầu tiên của Hikawa Kaigou trên màn ảnh ngoài các buổi phát sóng trực tiếp của các chàng trai, vào một ngày thứ Sáu rảnh rỗi.

Nắng chiều hắt nhẹ qua ô cửa kính. Trong căn phòng hai tầng phủ màu xanh thẫm, ba chàng trai cùng cô nữ Tổng quản đang mỗi người một việc riêng. Mỗi chàng trai là đại diện cho một thế hệ của nhà Sao.

\chrint

Ngồi nhàn nhã ở quầy bar, vị đàn ông đang nhai từng miếng pizza béo ngậy phảng phất hương vị thịt và phô mai chính là Arurandeisu (アルランディス, A-rư-ran-đê-i-xư) - mọi người thường gọi tắt cậu là ``Aruran''. Là một trong những thành viên đầu tiên thuộc thế hệ thứ Nhất của nhà \hls, người đàn ông gốc Ý này luôn mang đến một vẻ bí ẩn khó lường. Cậu từng tự nhận mình đã trải qua vô số cuộc đời với đủ vai trò: từ thám tử tư, tên trộm mảnh khảnh, giáo viên dạy học, cho đến cả một samurai lừng danh hay thậm chí một tay bắn tỉa cừ khôi, dù không ai có thể kiểm chứng được những câu chuyện đó, cậu vẫn luôn kể về chúng với ánh mắt tràn đầy nhiệt huyết.

Hiện tại, ước mơ lớn nhất Aruran theo đuổi là trở thành một thần tượng thực thụ của \hls. Cậu luôn đắm chìm sâu sắc vào cốt truyện của các nhân vật chính trong mọi tựa game mình chơi, đến nỗi nhiều lần người xem phải nhắc cậu trở lại thực tế. Điểm yếu duy nhất có lẽ là khả năng xác định phương hướng kém cỏi, khiến cậu hay bị lạc đường ngay cả ở những nơi đã đến nhiều lần, và tình yêu bất diệt với pizza, thịt và mọi loại nước ngọt có ga.

Trong mắt người hâm mộ cũng như các thành viên khác trong \hls, Aruran luôn là một người mang đến cảm giác ấm áp, vững chãi như một người cha. Chính sự chu đáo, chững chạc ấy đã giúp cậu có được những biệt danh thân thương như ``Papa'', ``Arupapa'' cùng nhiều tên gọi khác mà mọi người hay gọi đùa.

\model{25}{l}{7.5cm}{aruran1}

Không chỉ là một người anh cả đáng tin cậy, Aruran còn là một người thầy tuyệt vời, luôn sẵn sàng chia sẻ kiến thức của mình với mọi người. Cậu từng tổ chức nhiều buổi phát trực tiếp để giảng dạy về các chủ đề khoa học thú vị, từ quy trình tiêm chủng cho đến những thuật ngữ mạng phổ biến của Nhật Bản. Đồng thời, cậu cũng luôn cởi mở học hỏi những điều mới từ người khác, điều thể hiện rõ nhất qua những buổi stream cậu tập nói tiếng lóng tiếng Anh cùng các khán giả, khiến ai cũng bật cười vì sự dễ thương của cậu.

Vẻ ngoài trưởng thành của Aruran càng làm tôn lên vai trò ``người cha'' của mọi người: làn da ngăm rám nắng, mái tóc bạc ngắn gọn gàng cùng bộ râu dê gọn ghẽ, cặp mắt kính trên sống mũi và đôi mắt xám lúc nào cũng tràn đầy sự ấm áp. Cậu thường xuyên đeo một chiếc tai nghe màu đỏ cổ điển trên tai, tạo nên một nét cá tính riêng biệt không thể nhầm lẫn.

Trang phục thường ngày của cậu cũng toát lên sự chững chạc phù hợp với tính cách: chiếc áo khoác xanh lá cây có hai hàng cúc lịch sự, bên trong là áo thun đen đơn giản, phối cùng quần tây trắng thẳng thớm và cà vạt trắng có in họa tiết tinh tế. Cậu luôn đi đôi giày da đen bóng và không bao giờ quên đeo chiếc đồng hồ đeo tay thân thương trên cổ tay, món đồ mà cậu luôn nâng niu từ ngày đầu tiên gia nhập \hls.

Tiếp theo là Aragami Oga (\ruby{荒}{あら}\ruby{咬}{がみ}オウガ, \ruby{A-ra}{Hoang}-\ruby{ga-mi}{Giảo}\ Ô(-ư)-gaga) - chàng quỷ đang thong thả vươn vai, giãn gân cốt bên cạnh chiếc tủ lạnh mini màu trắng ngà ở góc phòng. Là thành viên thuộc thế hệ thứ Ba của nhà Ngôi Sao Vòm, Oga có một quá khứ không ai ngờ tới: một cựu binh đã trải qua năm nghìn năm chiến đấu trên chiến trường của một thế giới song song xa xôi. Sau khi giải ngũ và thảnh thơi du lịch đến Nhật Bản, số phận đã đưa đẩy cậu gặp Shinove - vị quản lý của \hls\ mà vai trò của ông tương đương với A-chan, dù hai người chưa bao giờ gặp mặt trực tiếp. Chính Shinove đã nhìn thấy tiềm năng ở Oga và săn đón cậu gia nhập tập thể \hls. Và từ khoảnh khắc được giới thiệu với thế giới game cổ điển, cậu như bị cuốn hút không thể dứt ra, mải mê chinh phục từng tựa game như thể chúng là những trận chiến mới cần chinh phục.

Ban đầu, những người bạn cùng nhóm TriNero của Oga từng không khỏi nhận xét rằng cậu là ``một người rất kỳ lạ'' khi mới gặp. Tuy nhiên, ẩn sau vẻ ngoài trầm mặc, khó đoán ấy lại là một tâm hồn vô cùng bình tĩnh và ôn hòa. Giọng nói của cậu lúc nào cũng trầm ấm, êm dịu, và cách cậu đối xử với người hâm mộ luôn tràn đầy sự ân cần, tạo nên một không khí ấm áp, thoải mái trong mọi buổi phát trực tiếp. Oga luôn trân trọng và vun đắp từng mối quan hệ với những người theo dõi mình, đến nỗi đã có lúc cậu thầm buồn khi nhận thấy sau buổi ra mắt, nhiều người đã chuyển từ cách xưng hô thân mật ``-kun'' (くん) sang cách gọi trang trọng hơn là ``-san'' (さん). Cảm giác đó như một khoảng cách nhỏ nhoi vô tình xuất hiện giữa cậu và những người mà mình luôn xem là bạn bè.

\model{25}{r}{7cm}{oga}

Vẻ ngoài gai góc của cậu cũng là một phần lý do ban đầu mọi người hơi ngại tiếp cận, nhưng càng chơi với cậu lâu, ai cũng nhận ra trái tim ấm áp mà Oga che giấu.

Về ngoại hình, Oga là một chàng quỷ cao lớn vượt trội, làn da ngăm rám nắng đi cùng vóc dáng săn chắc với cơ bụng từng đường nét rõ rệt. Điều khiến cậu nổi bật nhất là cơ ngực rắn rỏi, càng được tôn lên nhờ thiết kế trang phục độc đáo để lộ phần ngực, phô diễn nét mạnh mẽ của một cựu binh. Cậu sở hữu đôi mắt xanh lá cây trong trẻo, những chiếc răng nanh sắc bén đặc trưng của loài quỷ, mái tóc đen óng che một phần mắt trái, và đặc biệt là một chiếc sừng đơn lẻ có thể thay đổi cả màu sắc lẫn hình dạng theo tâm trạng. Chiếc sừng ấy chính là dấu ấn riêng biệt không thể nhầm lẫn của chàng quỷ Aragami Oga.

Trang phục thường ngày của cậu cũng phản ánh rõ nét cá tính mạnh mẽ nhưng vẫn có chút tỉ mỉ: một chiếc áo khoác đen được mở cúc, xăn tay lên lộ ra cổ áo cao màu xanh lá cây, phối cùng quần tây đen thanh lịch. Trên ngực trái là tấm che nhỏ in họa tiết hình con cọp xanh lá, càng làm tăng thêm nét hoang dã, mạnh mẽ. Cậu còn đeo một chiếc vòng cổ đính đá gắn liền với chiếc cà vạt hoa văn, đôi ủng cao cổ màu đen với phần đế và cổ ủng màu xanh lá rực rỡ. Trên hông, Oga đeo một băng đạn cũ chứa ba chiếc túi nhỏ màu nâu - theo lời cậu, bên trong có cả chiếc Game Boy Advance, chiếc máy chơi game cầm tay Neo Geo Pocket và vô số băng game cổ điển mà cậu nâng niu. Bên cạnh đó còn có một chiếc túi nhỏ khác đeo trên đùi, đôi găng tay lộ ngón, một chiếc khuyên bạc tròn đơn lẻ và miếng che mắt đen có thể tháo ra bất cứ lúc nào cậu muốn.

Cuối cùng, vị trí ở góc sô-pha thuộc về Yukoku Roberu (\ruby{夕}{ゆ}\ruby{刻}{こく}ロベル, \ruby{Giu}{Tịch}-\ruby{cô-cư}{Khắc}\ Rô-bê-rư), thành viên thế hệ thứ hai của nhà Ngôi Sao Vòm. Chính cậu là người quản lý quán bar kín đáo có tên ROBEL - một cái tên được đặt dựa trên chính tên riêng của mình. Roberu thường xuyên mở livestream vào những lúc quán chưa mở cửa, hay bất cứ khoảng thời gian rảnh rỗi nào trong ngày. Cậu nói chuyện với mọi người khá cởi mở, nhưng không ai có thể dễ dàng đoán được cậu thực sự đang nghĩ gì trong đầu. Dù cậu uống được khá nhiều rượu, thể chất lại yếu đến mức chỉ cần vài ly là đã có thể say xỉn. Khẩu hiệu mà Roberu luôn giữ trong đầu là: Con người sẽ luôn tin vào bất cứ điều gì họ muốn tin, bất kể sự thật có ra sao.

\model{26}{l}{7cm}{roberu}

Góc quầy bar nhỏ trong căn mancave chính là nơi Roberu thường lui tới nhất, nơi cậu có thể vừa pha chế đồ uống vừa trò chuyện thỏa thích với mọi người. Với cái miệng lanh lợi và nhanh nhảu của một chủ quán bar lâu năm, Roberu có thể nói liên tục về đủ loại chủ đề khác nhau. Có lúc cậu còn nói nhanh như máy bắn câu, liên miên gần một phút đồng hồ khiến người đối diện chẳng kịp trở tay. Hầu như trong mọi tình huống, cậu đều có thể lồng vào một câu đùa hài hước, và nhiều người thường khen ngợi cậu sở hữu khả năng làm một ``tsukkomi'' (người phản biện, đấm lố) xuất sắc. Chính vì lẽ đó mà Kaigou đã chọn cậu làm cầu nối trung gian giữa nàng Tổng quản và các thành viên còn lại của nhà Ngôi Sao Vòm. Thú vị hơn nữa, Roberu thường nói chuyện bằng phương ngữ Kansai - ngôn ngữ đặc trưng của vùng miền Tây Nhật Bản, càng làm tăng thêm nét duyên dáng và gần gũi của cậu.

Nếu quan sát kỹ, ngoại hình của Roberu cũng để lại ấn tượng khó phai trong lòng người gặp. Cậu mắc chứng dị sắc mắt hiếm gặp, với mắt phải mang màu cam ấm áp, trong khi mắt trái lại là màu tím huyền bí. Ngay dưới mắt trái còn có một nốt ruồi nhỏ, càng làm tăng thêm nét cá tính cho khuôn mặt. Mái tóc màu cam của cậu được chia đôi ở phần mái, xen lẫn những lọn tóc màu đen tương phản, cùng với chùm tóc nhỏ nhô lên - một chiếc ``ahoge'' đặc trưng không thể nhầm lẫn với ai khác.

Trang phục hàng ngày của Roberu cũng thể hiện rõ gu thẩm mỹ riêng của cậu. Bên trong là một chiếc áo thun xám với tay áo được xắn lên gọn gàng, bên ngoài khoác một chiếc vest nâu có túi nhỏ ở ngực, phối cùng quần tây đen thanh lịch. Cậu đi một đôi giày da đen thắt dây chữ thập, kết hợp với đôi tất màu cam nổi bật tạo điểm nhấn. Các phụ kiện khác còn có chiếc thắt lưng in họa tiết vằn ngựa vằn độc đáo, dải ruy băng cổ có thể chuyển màu, chiếc ghim cài áo hình sư tử mạnh mẽ, và một chiếc đồng hồ treo cổ cổ điển hoàn thiện bộ trang phục..

\chrint

Kaigou ngồi ở bàn ăn gỗ dài, sau vài ngày chung sống cùng ba chàng trai trong căn mancave này, nàng đã dần quen với không khí toàn nam giới, dù vẫn giữ chút đề phòng nhỏ trong lòng. Tay phải lướt nhẹ qua màn hình điện thoại, tay trái vừa rà soát lại dữ liệu công việc trên giao diện ảo chiếu từ chiếc đồng hồ đeo tay Kuon đã tặng mới - người cũng đã độ chế để Game có thể di chuyển vào đó hỗ trợ cô, cô khẽ thốt lên: ``Game, báo cáo lại số lượt xem của buổi stream tuần trước cho tôi xem nào.''

Giọng nói trên chiếc máy tính laptop lập tức nhấp nháy một dòng thông báo xanh lá trên màn hình ảo của đồng hồ Kaigou: ``Tổng lượt xem tuần trước đạt 1,2 triệu lượt, tăng 15\% so với tuần trước đó. Tương tác bình luận cũng tăng 22\%, chủ yếu là các bình luận tích cực về các màn chơi game của các thành viên.''

Cô Tổng gật gù hài lòng, ngón tay chạm nhẹ vào màn hình ảo của đồng hồ để mở thêm báo cáo doanh thu quảng cáo: ``Tốt lắm. Còn kế hoạch chuẩn bị cho buổi collab với nhóm \hls\ bên đó, các thủ tục đã xong hết chưa?''

``Chưa, còn thiếu chữ ký xác nhận từ phía quản lý của đội bên kia. Tôi đã gửi email nhắc lại lần thứ hai sáng nay, dự kiến sẽ nhận được phản hồi trong vòng 24 giờ tới,'' ông ta trả lời chính xác, các con số cập nhật liên tục trên màn hình trong suốt đoạn hội thoại.

Trong cuộc trò chuyện nhỏ đó, cậu chàng Chủ quán lại đang ngồi thẫn thờ trên chiếc sofa bọc vải màu xám đen. Cậu vò đầu bứt tóc, rồi cất lên một tiếng thở dài thượt: ``Chẳng lẽ\dots\ mình lại không có chút sức hút nào với mọi người sao\dots?''

\img{m1-1a}

Oga vươn người mở cửa chiếc tủ lạnh mini kêu két một tiếng, lấy ra một lon Coca-Cola mát lạnh. Nghe thấy tiếng thở ngắn thở dài của tiền bối, anh chàng lực điền ngoái đầu lại: ``Hửm? Tôi thấy anh cũng khá là cuốn hút đấy chứ, Robe-chan.''

Nàng Tổng ngẩng mặt lên khỏi màn hình điện thoại, tò mò nhìn cậu chàng tóc cam đang bất an: ``Xét về ngoại hình thì em trông cũng đâu đến nỗi nào đâu, Roberto.''

Nghe hai người lên tiếng khen ngợi, khuôn mặt rầu rĩ của Roberu lập tức rạng rỡ hẳn lên. Cậu nở một nụ cười toe toét, đắc ý vỗ đùi: ``Đúng chứ!? Quả nhiên là như vậy mà!''

Thế nhưng, chưa đầy hai giây sau, cậu chàng Quân nhân vừa bật nắp lon nước vừa bồi thêm một câu tĩnh rụi như tát nước vào mặt: ``Nhưng mà em thì còn `bánh cuốn' hơn anh nhiều.''

Kaigou không kìm được liền bật cười khẽ, lắc đầu: ``Phũ vãi phò, Oga à!''

Roberu nghe xong câu đó thì tức đến nghẹn họng. Cậu đứng xấc lên, chắp nắm đấm vào lòng bàn tay kêu rắc: ``Chú em có tin anh bẻ luôn cái sừng bên trái kia của chú không!?''

\img{m1-1b}

Tuy nhiên, cựu quân nhân Oga chẳng hề tỏ ra sợ sệt. Cậu điềm nhiên nhấp một ngụm nước ngọt rồi thách ngược lại: ``Cứ thử đi. Bẻ xong đảm bảo anh sẽ trở nên nổi tiếng nhất mạng xã hội luôn!''

Uất khuất vì không cãi lại được, Roberu quay sang nhìn cô nữ Tổng quản bằng ánh mắt cầu cứu. Nhưng Kaigou chỉ nhún vai lắc đầu ngao ngán: ``Thôi xin. Chị mà bẻ được cái sừng đó thì chị bẻ từ hôm đầu tiên rồi. Bẻ xong rồi bị một phát đuổi việc luôn thì chẳng đáng tí nào.''

Lúc này, Aruran đặt miếng pizza dở xuống dĩa, lấy khăn lau tay rồi thong thả đề xuất: ``Vậy thì sao em không thử tạo điểm nhấn bằng một chất giọng thật quyến rũ nhỉ?''

``Một giọng nói quyến rũ sao\dots?'' cậu chàng Chủ quán chớp mắt ngẫm nghĩ.

Kaigou gật gù: ``Ý kiến không tồi. Giọng em bình thường bắn liên thanh như máy khâu, lắm lúc kịch bản chị nghe còn không kịp dịch cơ.''

Cậu chàng người Ý mỉm cười: ``Một giọng nói trầm ấm luôn là vũ khí tuyệt vời đấy!''

Nghe tới đó, Oga đóng sầm cửa tủ lạnh lại, chậm rãi bước về phía góc phòng. Anh chàng nhàn nhã tựa một tay lên vách tường nhám màu tím, nghiêng đầu, hạ tông giọng trầm đục khêu gợi đầy lực hút: ``Nào\dots\ hãy thuộc về anh đi\~{}''

Bầu không khí trong căn phòng lập tức đóng băng. Kaigou, Roberu và Aruran nghếch cổ nhìn anh chàng lực điền đang vô cùng nghiêm túc\dots\ tán tỉnh góc tường. Mấy giọt mồ hôi hột to đùng chảy dài trên trán ba người.

\img{m1-1c}

Nàng Tổng khoanh tay, thở dài cất giọng: ``Này Oga. Chị ở đằng này cơ mà. Tường có nói chuyện được đâu.''

Cậu chàng Quân nhân khựng lại một nhịp, lờ đờ quay đầu sang: ``\dots À. Nhầm hướng mất rồi.''

Game từ mặt đồng hồ của Kaigou khẽ phát ra tiếng thông báo khô khốc: ``Phân tích nhịp tim Kaigou-user: Hoàn toàn ổn định ở mức 72 nhịp/phút. Màn thả thính của Aragami Oga đạt hiệu quả\dots\ 0\%.''

Chứng kiến màn trình diễn bá đạo của anh chàng cựu quân nhân, Roberu gật gù ngẫm nghĩ một lúc rồi đứng dậy kết luận: ``Anh nghĩ eanh bắt đầu hiểu tại sao người ta thích cái chất giọng trầm đó rồi\dots''

Oga vuốt cằm: ``Quả thực giọng trầm dễ thu hút người khác lắm.''

Aruran liếc nhìn cô gái duy nhất trong phòng rồi trêu: ``Cơ mà Hikawa-chan vẫn không suy suyển chút nào nhỉ?''

Kaigou bĩu môi đáp trả: ``Hai cậu nghĩ chị là fan cuồng mê muội giống như mấy con bánh bèo mê giai ngoài kia chắc? Mơ đi Diễm!''

Cậu chàng người Ý nhếch môi cười, quay sang vỗ vai Roberu: ``Dù sao thì\dots\ Roberu này, sao chú em không thử hạ cái giọng của mình xuống xem?''

Thế là cậu Chủ quán rượu bắt đầu thử nghiệm. Roberu đứng thẳng lưng, giơ tay chào gần giống phong cách nhà binh, cố hạ giọng xuống chút xíu: ``Xin chào, mình là Yukoku Roberu đây.''

Kaigou khẽ lắc đầu, không tiếc lời đánh giá: ``Nghe vẫn `trẻ trâu' lắm em.''

Oga gật đầu đồng tình: ``Ừ. Chưa tới đâu.''

Cậu chàng Thế hệ thứ Hai sững sờ một chút, rồi quyết định hạ giọng xuống thêm một nấc nữa. Cậu đặt tay lên ngực, hơi cúi người tỏ vẻ lịch thiệp: ``Xin chào, mình là Yukoku Roberu đây.''

Aruran vẫn chưa hài lòng: ``Chắc phải trầm hơn nữa mới đạt.''

Cậu chủ quán quyết định gồng mình hạ giọng xuống hết mức có thể, bỏ luôn cả động tác tạo dáng: ``Xin chào, mình là Yukoku Roberu đây\dots''

Cậu chàng Quân nhân xòe tay ra hiệu kéo xuống: ``Hạ thêm tí nữa!''

Roberu gồng cổ, mặt mày nhăn nhó như rặn đẻ, phát ra một âm thanh ồm ồm đục ngầu: ``Xin\dots\ chào\dots\ mình\dots\ là\dots\ Yu\dots\ ko\dots\ ku\dots\ Ro\dots\ be\dots\ ru\dots''

Kaigou toát mồ hôi hột, lùi lại nửa bước: ``Trời đất ơi\dots\ Em nghe giống mấy ông cụ tám mươi tuổi bị viêm họng quá\dots''

Aruran lại tỏ ra hứng thú: ``Này, nghe có vẻ triển vọng đấy chứ!''

Nàng Tổng thì nhăn mặt: ``Chị thì hoàn toàn không nghĩ thế\dots''

Cậu chàng người Ý rút điện thoại ra, mở ứng dụng trợ lý ảo rồi hướng micro về phía Roberu: ``Nói lại lần nữa vào đây xem nào.''

\img{m1-1d}

Cậu chủ quán nhắm mắt ráng gồng cái giọng cao niên đó lần nữa: ``Xin chào, mình là Yukoku Roberu đây\dots''

Cả phòng im lặng chờ đợi. Tưởng rằng chiếc máy sẽ không thể nhận ra giọng nói, nhưng rồi màn hình sáng lên, giọng điện tử quen thuộc của Game cất lên đầy xóc xỉa: ``Xin chào, cụ Yukoku. Cần tôi đặt xe lăn điện cho cụ để ngày mai đi khám bác sĩ tai mũi họng không?''

Cậu chàng Chủ quán sững sờ gục ngã: ``Trời ơi! Không được rồi!''

Oga và Aruran chỉ biết cười trừ, trong khi Kaigou không kìm nổi nữa, bật cười phá lên: ``Haha! Game biết đùa dễ sợ! Giọng ồm như cưa gỗ thế kia người ta lại tưởng em là ông già không bằng!''

``Lạ nhỉ\dots'' Oga gãi đầu.

Mặt Roberu đỏ bừng vì cay đắng. Cậu hất tay gạt phăng ý định hỗ trợ của mọi người: ``Thôi! Méo để mọi người giúp nữa! Một mình tôi tự làm!''

Nàng Tổng cố nhịn cười, xua tay hòa giải: ``Ấy ấy, đừng dỗi mà `ông chủ'! Chị xin lỗi, chị xin lỗi mà!''

Oga tặc lưỡi: ``Có vẻ như anh không thể cuốn hút mọi người bằng cái giọng cưa gỗ đó rồi nhỉ\dots''

Kaigou lườm cậu ta một cái: ``Nào Oga. Đừng có trêu em ấy nữa.'' nhưng miệng thì cố nhịn tiếng cười lại.

``Bố già'' Aruran tiến lại gần, nhẹ nhàng đặt tay lên vai anh chàng đang phừng phừng tự ái: ``Thế nếu không dùng cách đó thì chú em định làm gì tiếp theo?''

Roberu khựng lại, nhận ra câu hỏi của Aruran hoàn toàn có lý. Cậu đứng ngập ngừng giữa phòng. Thấy vậy, Kaigou và Oga cùng bước đến bên cạnh, chân thành động viên cậu chàng đang bất lực.

Cậu chàng Quân nhân bỗng nhiên thay đổi thái độ, nét mặt trở nên nghiêm túc đột xuất: ``Anh không phải lo đâu. Lần này bọn em sẽ thực sự nghiêm túc giúp anh.''

Kaigou mím môi mỉm cười: ``Trêu thế thôi, chứ chị nhìn em có vẻ rất tâm huyết với việc xây dựng hình ảnh mà.''

Roberu ngoái đầu lại, bán tín bán nghi hỏi: ``Có đúng là lần này giúp thật không đấy\dots?''

Chưa kịp để cậu định hình, Oga bất ngờ tiến sát lại, ngọn lửa quyết tâm của một cựu quân nhân bùng cháy rừng rực quanh người. Anh trừng mắt nhìn đàn anh bằng một thần thái cực kỳ nghiêm khắc: ``Nhưng anh phải nghe theo tất cả những gì bọn em bảo! \textbf{RÕ CHƯA!?}''

Áp lực khí thế quân đội từ cậu chàng bùng lên dữ dội đến mức khiến Kaigou - dù vốn dĩ hổ báo và khắt khe - cũng phải toát mồ hôi hột, khẽ giơ tay can ngăn: ``Em không cần phải ngặt nghèo với em ấy như đi nghĩa vụ quân sự thế đâu Oga\dots''

Trái ngược hoàn toàn với sự lo lắng của cô Tổng, Roberu lại như được tiếp thêm năng lượng. Cậu lập tức đứng thẳng lưng tắp, giơ tay chào chuẩn phong cách nhà binh, gào lên đầy quyết tâm: ``\textbf{Sơ! Dét, sơ!} (Rõ, thưa sếp!)''

\img{m1-1e}

Cô nàng Tổng quản ngơ ngác mất vài giây rồi chỉ biết bất lực cười trừ cho qua. Oga gật đầu đầy hài lòng: ``Tốt. Chúng ta cùng bắt đầu buổi tập nào!''

Kaigou khoanh tay, mỉm cười khích lệ: ``Cố gắng lên nhé, `thanh niên lắm mồm'!''

\ct

Oga đập mạnh hai tay xuống bàn, đứng phắt dậy với tác phong như một vị tướng quân chuẩn bị phát lệnh: ``Bài tập Nổi tiếng số một: Lấy bằng cấp!''

Aruran nhanh tay rút từ trong túi áo ra một tờ phiếu quảng cáo gấp ba nham nhở, nhè nhẹ nhét vào tay cậu chàng tóc cam: ``Đi học lấy chứng chỉ chuyên viên phục vụ khăn tắm cấp độ 3 đi nào, chàng trai!''

Roberu trợn tròn mắt, nhìn tờ giấy kẹp trong tay như thể nhìn một sinh vật lạ ngoài hành tinh: ``Cái quái gì đây!? Tại sao một streamer như em lại phải đi học lấy bằng phục vụ khăn tắm chứ!?''

Mặt đồng hồ của Kaigou rung nhẹ, giọng Game cất lên thong thả và đầy tính sát thương: ``Trích xuất dữ liệu: Tỷ lệ khán giả đăng ký kênh tăng lên nhờ chứng chỉ gấp khăn tắm ước tính đạt khoảng 0,0001%. Một khoản đầu tư vô cùng chiến lược đấy, Roberu-user.''

``\textbf{Rõ ràng là mỉa mai còn gì nữa!!}'' Roberu gào lên.

Không để hậu bối kịp hạ hỏa, cựu quân nhân Oga đã tiếp tục chuyển sang bài học số hai. Cậu khoanh hai tay trước ngực, giọng trầm đục đầy sức nặng: ``Bài tập số hai: Rèn luyện tính chịu đựng và sự hung hãn!''

Aruran cúi xuống góc ghế bành, lôi ra một cây gậy bóng chày bằng gỗ rồi chĩa về phía Roberu: ``Cầm lấy cây này, ra ngoài hành lang đập nát mấy cái cục nóng điều hòa đi!''

Kaigou đang lướt điện thoại suýt nữa sặc nước bọt. Cô ngẩng đầu nhìn ``Bố già'' người Ý bằng ánh mắt kỳ quặc: ``Ơ này Bố già\dots\ Đập nát điều hòa thì tới mùa hè cả cái nhà kho này chịu chết nóng à? Với lại đây là tài sản công ty đấy!''

Roberu ôm đầu bất lực: ``Lố bịch! Quá sức lố bịch luôn ấy!'' Đáp lại cậu chỉ là một nụ cười trầm ấm nhưng gian xảo tới từ vị trí của anh chàng người Ý ấy.

Oga vuốt cằm, nhíu mày ngẫm nghĩ về bài tập tiếp theo. Nhưng dường như chất xám của anh chàng lực điền đã cạn kiệt, Oga tặc lưỡi nói đại: ``Thế thì\dots\ đi mua cho bọn em ít bánh mì nóng về đây!''

Aruran chêm vào ngay lập tức: ``Nhớ là phải vừa mới ra lò, vỏ còn giòn rụm và ấm nóng đấy nhé!''

Roberu cay đắng chỉ tay vào mặt mình: ``Thế rốt cuộc em là streamer tìm kiếm sự nổi tiếng hay là chân sai vặt cho hai anh đấy!?''

Kaigou lắc đầu, rút từ trong ví ra vài tờ tiền rồi ấn vào tay cậu chủ quán: ``Này Oga, Aruran, hai cậu có chân có lực thì tự đi mà mua chứ. Đừng có bóc lột em nó thế\dots\ Cơ mà thôi, Roberto này, cầm tiền đi mua giùm chị ổ mì pa-tê luôn nhé.''

Dù bực tức đến mức đầu muốn bóc khói, cậu chủ quán vẫn hậm hực cầm lấy xấp tiền của cô Tổng rồi dậm chân thình thịch đi ra ngoài.

\begin{figure}[H]
  \centering
  \begin{subfigure}[t]{0.45\textwidth}
    \centering
    \includegraphics[width=8cm]{../../../../Image/BookM1/m1-1f1.png}
  \end{subfigure}
  \quad
  \begin{subfigure}[t]{0.45\textwidth}
    \centering
    \includegraphics[width=8cm]{../../../../Image/BookM1/m1-1f2.png}
  \end{subfigure}
  \quad
  \begin{subfigure}[t]{0.45\textwidth}
    \centering
    \includegraphics[width=8cm]{../../../../Image/BookM1/m1-1f3.png}
  \end{subfigure}
  \caption{Những bài tập Nổi tiếng mà Oga đề ra cho Roberu, nhưng lại biến thành trò cười.}
\end{figure}

Nhìn cánh cửa đóng sầm lại, Kaigou khoanh tay, nheo mắt nhìn hai chàng trai đang ngồi nhàn nhã ở quầy bar: ``Chị có cảm giác hai người đang đem em ấy ra làm trò hề để tiêu khiển thì đúng hơn đấy.''

Oga gãi gãi cái sừng bên trái, tỉnh rụi đáp: ``Thế á? Tại tự dưng em thấy đói bụng quá\dots''

Aruran thở dài thong thả: ``Có vẻ như chiến thuật vừa rồi không hiệu quả ha.''

``Theo thống kê của tôi,'' Game cêm vào, ``hành vi vừa rồi của hai cậu được xếp vào loại `bóc lột lao động trẻ em cấp độ nhẹ'. Mặc dù Yukoku-san đã 28 tuổi.''

Kaigou xoa thái dương: ``Mấy cái trò ba lăng nhăng đó không giúp em ấy nổi tiếng hơn đâu. Nghĩ cách nào đàng hoàng chút đi!''

\ct

Mười phút sau, Roberu xách một túi bánh mì nóng hổi bước vào. Cậu đập mạnh túi bánh xuống bàn ăn, rồi hầm hực lao đến chiếc sô-pha bọc vải xám, đổ ập người xuống ngồi giận dỗi. Áp lực uất khuất tỏa ra xung quanh cậu to lớn đến mức Kaigou đành phải thở dài, tiến lại ngồi xuống bên cạnh nhẹ nhàng vỗ lưng an ủi: ``Thôi nào, làm tốt lắm Roberto. Ít ra em cũng mua đúng loại bánh giòn rụm rồi.''

Ở bàn ăn, Oga và Aruran vừa nhâm nhi từng miếng bánh mì thơm phức vừa vò đầu bứt tóc tìm giải pháp thực sự. Cậu chàng Quân nhân vừa nhai vừa lẩm nhẩm: ``Nếu như mấy cái đó không được\dots\ thì rốt cuộc anh ấy đang thiếu cái gì nhỉ? Một cái gì đó\dots\ tinh tế hơn\dots''

Cậu chàng người Ý búng tay *TÁCH* một cái: ``Sự nhã nhặn! Sự thanh lịch của một quý ông!''

``Chính nó!'' Oga quay phắt đầu sang phía sofa, gào lên, ``\textbf{Này Robe-chan! Anh thử làm theo cách này xem nào!}''

\ct

Ít phút sau, buổi thực hành ``Sự nhã nhặn'' chính thức bắt đầu.

Một bản nhạc vĩ cầm cổ điển du dương (từ Game phát ra) vang lên trong căn phòng. Roberu ngồi thẳng lưng trên ghế đẩu quầy bar, ngón tay út vẫy nhẹ khi nâng tách trà gốm sứ lên nhấp một ngụm đầy phong thái hoàng gia.

Cậu nhẹ nhàng đặt tách trà xuống đĩa, mỉm cười tao nhã, ánh mắt chiếu thẳng vào ống kính rồi cất giọng bằng chất giọng đài các nhất có thể: ``Xin phép mọi người\dots\ cho tôi đi ra ngoài giải quyết nỗi buồn một chút nhé\~{}''

\img{m1-1g}

Nói rồi, cậu chống hai tay lên hông, nở một nụ cười cực kỳ\dots\ hớ hênh và khó đỡ.

*XẸT!*

Bản nhạc vĩ cầm phía sau đột ngột bị cắt đứt phũ phàng. Màn biểu diễn khép lại, khuôn mặt cậu chàng Chủ quán lập tức nghệt ra, trở lại vẻ ngơ ngác thường ngày: ``\dots Kiểu như thế này đúng không hai người?''

Oga và Aruran nhìn nhau chớp chớp mắt, hoàn toàn câm nín. Cả hai đồng loạt quay sang nhìn cô nữ Tổng quản để xin đánh giá chuyên môn.

Kaigou ấp úng, khóe miệng giật giật, không biết phải dùng từ ngữ nào để diễn tả những câu nói ngớ ngẩn vừa rồi: ``Ờm\dots\ Chị thì không nghĩ cái việc thông báo đi vệ sinh một cách điệu đàng như thế được gọi là `nhã nhặn' đâu em\dots''

Game phát loa cất giọng khô khốc: ``Cậu đang nhầm lãn tai hại giữa phong thái quý tộc nước Anh và hành vi sinh lý cơ bản của con người rồi. Khán giả sẽ bỏ chạy trong vòng 3 giây cho mà xem.''

Roberu uất khuất nhảy dựng lên: ``\textbf{Chính hai người bảo tôi phải làm ra vẻ thanh lịch tao nhã cơ mà!}''

Cậu chàng Quân nhân giơ hai tay hòa hoãn: ``Này, em đang cố giúp anh đấy nhé. Bình tĩnh lại đi nào.''

Cậu người Ý tặc lưỡi: ``Thôi nào, ta cùng suy nghĩ cách khác vậy.''

Kaigou tựa lưng vào sofa, thì thầm với Game: ``Chị bắt đầu nghi ngờ về năng lực cố vấn của hai ông tướng này rồi đấy\dots''

``Để cho mấy diễn viên hài góp ý thì đúng là hết cứu thật rồi,'' giọng nói ở trong chiếc đồng hồ thở dài đáp lại.

\ct

Một lúc sau, Aruran chắp tay sau lưng, đi đi lại lại quanh chiếc sofa như một triết gia đang tìm kiếm chân lý. Kaigou và Roberu ngồi gục gù trên ghế nhìn theo hướng chuyển động của vị ``Bố già''.

Kaigou khoanh tay hỏi: ``Nghĩ được ý tưởng gì mới chưa, Bố già?''

Cậu chàng dừng bước, ngước mắt nhìn trần nhà đầy đăm chiêu: ``Ta hãy thử thay đổi cách tiếp cận vấn đề\dots''

Roberu thở dài thượt: ``Anh nói đi\dots\ em nghe đây\dots''

``Em nghĩ kiểu người đàn ông nào mới thực sự là một người nổi tiếng và có sức hút?'' Aruran đặt câu hỏi rồi tự mình đưa ra câu trả lời bằng một chất giọng đầy triết lý, ``Đó là khi người đàn ông đó\dots\ thực sự tin tưởng vào chính bản thân mình!''

Lời nói của Aruran như một luồng sáng chói lọi rọi thẳng vào tâm trí của hậu bói. Như được khai sáng, cậu Chủ quán ngẩng phắt đầu lên, đôi mắt rực cháy niềm hy vọng: ``X-Xin phép mọi người em ra ngoài một chút!!''

Chưa kịp để ai phản ứng, Roberu đã vắt chân lên cổ lao thẳng ra khỏi trụ sở.

Aruran thở dài gãi đầu: ``Chẳng biết chú em có thực sự hiểu đúng ý của em không nữa\dots''

Kaigou cười gượng gạo, vỗ vai ``Bố già'': ``C-Chị nghĩ là\dots\ chắc mọi thứ sẽ ổn thôi\dots\ hy vọng thế.''

\ct

Tối hôm đó, trên kênh phát sóng trực tiếp của Yukoku Roberu xuất hiện một buổi lên sóng đặc biệt với tiêu đề: ``Một cái tôi mới''.

Trên màn hình, cậu ngồi thẳng lưng tắp, mặc trang phục chỉn chu, biểu cảm nghiêm nghị đến mức cứng đờ. Cậu cất giọng bằng một sự nghiêm túc kỳ lạ, hoàn toàn đánh mất sự náo nhiệt, nhí nhảnh thường ngày: ``Buổi sáng tốt lành mọi người. Chúng ta hãy cùng dành cho nhau khoảng thời gian tuyệt đẹp hôm nay nhé\dots''

Tại căn mancave, bốn người bao gồm Kaigou, Oga, Aruran và Roberu đang cùng nhau xem lại đoạn video phát lại và lướt phần bình luận của khán giả.

\img{m1-1h}

Roberu trợn tròn mắt, mặt cắt không còn một giọt máu: ``Em\dots\ em không ngờ phản ứng của mọi người lại thành ra thế này\dots''

Aruran chỉ tay vào màn hình: ``Khán giả còn đang spam toàn chữ `草' (lol/cỏ) kín cả khung chat kìa!''

``''\textbf{Hahahahaha! Trời ơi, chết mất thôi!}''

Kaigou không kìm chế được nữa, cô ôm bụng ngã rần xuống sàn nhà, lăn lộn cười ngặt nghẽo đến mức chảy cả nước mắt. ``Cười chết chị mất!! Trông em nghiêm túc quá mức luôn ấy Roberto!! Nhìn giống hệt một con NPC đang cố bắt chước làm người thật! Haha!''

``Cảnh báo,'' Game cất giọng cảnh báo từ trên màn hình máy tính, ``lLượng oxy tiêu thụ của Kaigou-user đang giảm nhanh do cười quá sức. Nếu tiếp tục lăn lộn thêm 30 giây, cô ấy sẽ đạt vận tốc xoay của cánh quạt trực thăng.''

Oga khoanh tay nhìn cô Tổng đang cuộn tròn trên sàn: ``Cứ cái đà này có khi chị ấy biến thành cái trực thăng cất cánh thật chứ chẳng chơi.''

Aruran tặc lưỡi trêu: ``Hoặc là chúng ta sẽ có cả một bãi cỏ xanh mướt trong phòng nhờ đống chữ `草' kia\footnote{Trong tiếng Nhật, ``草'' (đọc là kusa) đồng nghĩa với ``haha'', ``lol'' hay ``XD'' thường dùng ở Việt Nam và các nước trên thế giới, có nghĩa là cười. Nó bắt nguồn từ ``wwwwwwwww'' (đọc là warau), nhìn giống cỏ. Ý của Oga là buổi phát sóng đó làm cho fan của Roberu ôm bụng cười, spam ``wwwww'' và ``草'' liên tục, làm cho người ta cảm tưởng rằng họ sẽ có một bãi cỏ mọc um sùm ở đó.} đấy.''

Roberu méo xệch cơ mặt, uất khuất nhìn nàng Tổng quản đang rên hừ hừ vì cười quá đau bụng: ``\dots K-Kaigou-senpai\dots\ chị\dots\ chị cười xong chưa vậy\dots?''

Phải mất một lúc lâu sau, Kaigou mới lồm ngồm bò dậy, lau vệt nước mắt nơi khóe mắt rồi cùng Oga và Aruran tiến lại gần anh chàng tóc cam đang gục đầu thất vọng trên bàn.

Aruran nhẹ nhàng vỗ vỗ vào vai Roberu, đúc kết bằng sự từng trải: ``Con đường trở thành người nổi tiếng quả là gian nan và chông gai thật đấy!''

Kaigou cũng nhếch môi mỉm cười, đặt tay lên bờ vai bên kia của cậu chủ quán, vỗ nhẹ động viên: ``Cũng phải ha\dots\ `Thanh niên lắm mồm' của chúng ta còn phải cố gắng nhiều lắm.''

Và đó chính là những gì đã diễn ra trong ngày làm việc chính thức đầu tiên của Hikawa Kaigou tại quán rượu Roberu. Một ngày không có quá nhiều biến cố lớn, chỉ tràn ngập những tiếng cười và sự náo nhiệt kỳ lạ giữa cô Tổng quản cùng ba chàng trai.

Hoặc có thể\dots\ sự bình yên này chỉ là khoảng lặng trước những cơn bão hỗn loạn tiếp theo đang chờ đợi họ ở phía trước. Tương lai mà, ai biết trước được điều gì chứ?

%==========================================TRUYỆN PHỤ=========================================
\omk

Sau khi sóng gió tạm qua đi, cả bốn người lại quây quần quanh bàn ăn để dọn dẹp đống vỏ bánh mì. Roberu kéo ghế ngồi xuống, chống cằm tò mò lên tiếng: ``Vậy\dots\ rốt cuộc kiểu người đàn ông mà chị cho là có sức hút thực sự là như thế nào vậy, Hikawa-san?''

Oga và Aruran cũng lập tức hướng ánh mắt về phía cô nữ Tổng quản. Anh chàng cựu quân nhân gật gù: ``Đúng đấy. Chị chưa nói cho bọn em biết gu của chị là gì đâu nhé.''

Cậu chàng người Ý mỉm cười chêm vào: ``Nhắc mới nhớ, tụi em còn chẳng biết chị đã có bạn trai hay chưa nữa đấy.''

Nghe thấy thế, Kaigou bật cười tủm tỉm, ngả người tựa vào lưng ghế: ``Đương nhiên là chị chưa có rồi. Nhưng mà\dots\ chị có một người mà chị đang rất thích đấy.''

Cả ba chàng trai đồng loạt rướn người về phía trước, ánh mắt sáng rực sự hóng hớt: ``Ai thế chị!?'' Nàng Tổng chần chừ một chút, rồi cố tỏ ra khách quan nhất có thể: ``Đó là một người mà chị mới quen gần đây thôi. Nói chung là\dots\ anh ấy có khá nhiều điểm giống chị.''

Ba người nghe vậy liền ``Ồ\~{}'' lên một tiếng đầy thấu hiểu, gật gù liên tạch.

Kaigou giơ một ngón tay lên đếm: ``Kiểu như\dots\ anh ấy cực kỳ nghiêm túc trong công việc này\dots''

Oga nhếch môi, phũ phàng cắt ngang: ``Nãy giờ em toàn thấy chị bày trò chọc phá Robe-chan thì có.''

Kaigou phồng má, lườm anh chàng cựu quân nhân một cái cháy mắt: ``Mấy đứa không thấy đấy thôi! Chứ những lúc chị bắt tay vào làm việc là nghiêm chỉnh cực kỳ luôn nhé!''

Roberu gật gù đứng ra gỡ gạc giúp cô Tổng: ``Cái này thì em công nhận đúng nè. Hôm vừa rồi chị ấy hỗ trợ em lắp đặt và trang trí cái quầy bar này hăng say lắm, đến mức quá giờ ăn trưa mà tay vẫn không ngơi nghỉ chút nào.''

Aruran ngạc nhiên nhìn nàng Tổng: ``Ồ, vậy là chị cũng có những khoảnh khắc năng nổ như thế thật à?''

``Có chứ sao không!'' Kaigou khoanh tay bĩu môi. ``Mấy đứa nghĩ chị lười nhác lắm chắc?'' Ba người con trai nhìn nhau, rồi đồng thanh trả lời không chút do dự: ``Thì đúng thế còn gì nữa.''

Mặt đồng hồ của Kaigou bất ngờ sáng lên, giọng nói đều đều vô cảm của Game vang lên phá tan bầu không khí: ``Phân tích nhật ký hoạt động 7 ngày qua của Kaigou-user: 68\% thời gian dùng để lướt điện thoại và chơi game, 20\% dành cho việc ăn uống, và chỉ 12\% là thực sự làm công việc của một Tổng quản. Kuon-user vẫn chưa biết chuyện này đâu.''

Kaigou toát mồ hôi hột, lén lút lấy tay che chiếc đồng hồ lại, mặt đỏ bừng vì bị ``bóc phốt'' trắng trợn.

``\textit{Chết thật\dots\ Mình phải thực sự làm gương cho lũ đàn em này mới được\dots}'' cô nghĩ thầm trong đầu, rồi hắng giọng một cái thật to để chữa thẹn: ``Khửm! Bỏ qua chuyện đó đi! Ngoài sự nghiêm túc ra thì\dots\ anh ấy còn cực kỳ tận tình và tốt bụng nữa. Giống hệt chị luôn!''

Oga gật đầu thừa nhận: ``Cái này thì em phải công nhận. Chính chị đã giúp bọn em lắp đặt bộ máy tính livestream với sửa lỗi phần mềm mấy hôm trước.''

Aruran hưởng ứng: ``Công nhận tay nghề kỹ thuật của chị tốt thật.''

Roberu rướn người hỏi tiếp: ``Thế ngoài ra anh ấy còn điểm gì nữa không chị?''

``À, còn nhiều lắm chứ\dots''

Kaigou bắt đầu hào hứng kể lại những ưu điểm của người đồng nghiệp đang ngày đêm ``chinh chiến'' bên nhánh \hll. Sự chân thành và ánh mắt lấp lánh của cô khi nhắc về người đó khiến ba chàng trai \hls\ dần hiểu ra kiểu người đàn ông có thể đốn gục trái tim của cô Tổng quản.

Cậu chàng người Ý thở dài thong thả: ``Tìm được một người hoàn hảo như thế thời nay đúng là khó thật đấy.''

Cậu chàng Quân nhân gật gù: ``Công nhận.''

Nàng Tổng lắc đầu cảm thán: ``Đàn ông bây giờ có nhiều chủng loại lắm. Tốt có, xấu có, ẩm IC cũng có\dots''

Roberu nghiêng đầu, đưa ra một câu hỏi mang tính quyết định: ``Nhưng vấn đề là\dots\ anh ấy có đáp lại tình cảm của chị không đấy?''

Không một giây chần chừ, Kaigou ưỡn ngực quả quyết: ``Chị tin chắc là có chứ!'' Mặc dù sâu thẫm trong lòng, chính cô cũng chưa từng nghĩ đến chuyện tiến xa hơn với người đã dang tay đón nhận và cho cô một gia đình mới.

Mặt đồng hồ lại rung nhẹ, Game phát loa nhỏ đủ để một mình cô nghe thấy: ``\hl{Đừng tự tin thái quá như thế. Cậu ta cũng chỉ coi cô là một cô em gái hơn là một cô bạn gái đấy.}''

Nàng Tổng cốc nhẹ vào mặt đồng hồ một cái cạch để dập tắt giọng của chiếc trợ lý ảo lắm chuyện.

Dù tỏ ra khá lo ngại cho sự tự tin bộc trực của cô nàng, ba chàng trai vẫn mỉm cười chúc phúc cho mối quan hệ của cô. Kaigou xua tay, quay trở lại chủ đề chính: ``Mà chị cũng đâu có bắt các em phải trở thành mẫu người như anh ấy đâu. Các em cứ là chính mình là tốt nhất rồi.''

Roberu ngơ ngác: ``Ý chị là\dots''

Aruran mỉm cười giải thích thay: ``Tức là làm những gì mà mình giỏi nhất. Anh nghĩ ý chị Kaigou là như thế.''

Oga vuốt cằm: ``Em đoán đó mới chính là cách thực sự để một người đàn ông tạo nên sức hút riêng cho bản thân.''

``Chính xác!'' Kaigou vỗ tay bộp một cái, ``Đặc biệt là với nghề idol như các em, mỗi người nên có một phong cách riêng biệt. Chứ ai cũng giống ai thì khán giả xem chán chết.''

Những lời chốt hạ của cô Tổng như khai thông hoàn toàn bế tắc trong lòng cậu chủ quán. Roberu sáng mắt lên, gập người cảm ơn đầy chân thành: ``Em đã vỡ ra được nhiều điều rồi\dots\ Cảm ơn chị rất nhiều, Kaigou-senpai!''

Nàng Tổng mỉm cười xua tay: ``Không cần khách sáo đâu. Mà này, từ giờ cứ gọi chị là Kaigou được rồi. Chị đây cũng đang phải tập cách nói chuyện và tiếp xúc với đám con trai bọn em nhiều hơn nữa đây.''

Aruran tò mò chớp mắt: ``Ơ\dots\ Chị có thù hằn gì với con trai bọn em sao?''

Kaigou im lặng một nhịp. Ánh mắt cô khẽ đờ ra, khóe miệng nở một nụ cười gượng gạo nhưng đôi mắt lại lóe lên một tia sáng nguy hiểm đến rợn người: ``\dots Cứ coi như là vậy đi. Cầu cho chị không nổi điên lên mà rút cái nơ trên đầu để đánh các em là tốt lắm rồi.''

Khí áp trong phòng đột ngột giảm sâu. Roberu rén đến mức thụt cổ lại, còn Oga và Aruran - những người vốn dày dặn kinh nghiệm sống - cũng lập tức nâng cao cảnh giác, tự dặn lòng không nên chọc giận cô Tổng quản này dưới bất kỳ hình thức nào.

Tuy nhiên, vượt lên trên những điểm bất đồng và sự e dè ban đầu, mục tiêu chung của cả bốn người lúc này đã trở nên rõ ràng hơn bao giờ hết: cùng nhau đoàn kết, nỗ lực hết mình để đưa cái tên \hls\ vươn xa, và quan trọng nhất là hoàn thành ước mơ cháy bỏng của vị CEO kính yêu - YAGOO.

Mùa giải của những ``idol nam kiêm diễn viên hài'' \hls\ chính thức bắt đầu!

\end{document}